\section{Diagonalization of a Symmetric Linear Map} 
\bx{
\en{
    \item 
    \item 
    \item 
    Let $V$ be a finite dimensional vector space with a positive definite scalar product. 
    Let $A: V \to V$ be a symmetric linear map. 
    We say that $A$ is \textbf{positive definite} if $\langle Av, v \rangle > 0$ for all $v \in V$ and $v \neq O$. 
    Prove: 
    \ea{
        \item 
        If $A$ is positive definite, then all eigenvalues are $> 0$. 
        \item 
        If $A$ is positive definite, then there exists a symmetric linear map $B$ such that $B^2 = A$ and $BA = AB$. 
        What are the eigenvalues of $B$? 
        [\textit{Hint}: Use a basis of $V$ consisting of eigenvectors.] 
    }
    \begin{Solution}
        \ea{
            \item 
            \begin{proof}
                Let $v$ be an eigenvector $\neq O$ with eigenvalue $\lambda$. 
                Then $\langle Av, v \rangle = \langle \lambda v, v \rangle = \lambda \langle v, v \rangle$. 
                Since $\langle v, v \rangle > 0$ and $\langle Av, v \rangle > 0$, it follows that $\lambda > 0$. 
            \end{proof}
            \item 
            \begin{proof}
                Pick a basis of $V$ consisting of eigenvectors. 
                The vector space $V$ can then be identified as the space of coordinate vectors with respect to this basis. 
                The matrix of $A$ then is a diagonal matrix, whose diagonal elements are the eigenvalues, and are therefore positive. 
                We can then let $B$ have diagonal elements equal to the square roots of those of $A$. 
            \end{proof}
        }
    \end{Solution}

    \item 
    We say that $A$ is \textbf{positive semidefinite} if $\langle Av, v \rangle \geq 0$ for all $v \in V$. 
    Prove the analogues of (a), (b) of Exercise 3 when $A$ is only assumed semidefinite. 
    Thus the eigenvalues are $\geq 0$, and there exists a symmetric linear map $B$ such that $B^2 = A$. 
    \begin{solution}
        \ea{
            \item 
            \begin{proof}
                Let $v$ be an eigenvector $\neq O$ with eigenvalue $\lambda$. 
                Then $\langle Av, v \rangle = \langle \lambda v, v \rangle = \lambda \langle v, v \rangle$. 
                Since $\langle v, v \rangle \geq 0$ and $\langle Av, v \rangle \geq 0$, it follows that $\lambda \geq 0$. 
            \end{proof}
            \item 
            \begin{proof}
                Pick a basis of $V$ consisting of eigenvectors. 
                The vector space $V$ can then be identified as the space of coordinate vectors with respect to this basis. 
                The matrix of $A$ then is a diagonal matrix, whose diagonal elements are the eigenvalues, and are therefore nonnegative. 
                We can then let $B$ have diagonal elements equal to the square roots of those of $A$. 
            \end{proof}
        }
    \end{solution}
    \item 
    Assume that $A$ is symmetric positive definite. 
    Show that $A^2$ and $A^{-1}$ are symmetric positive definite. 
    \begin{Solution}
        \begin{proof}
            From ${^{t\!}(AA)} = {^{t\!}A}{^{t\!}A} = AA$, it follows that $A^2$ is symmetric. 
            Furthermore, for $v \neq O$, 
            \begin{equation*}
                \langle A^2 v, v \rangle = \langle Av, {^{t\!}A} v \rangle = \langle Av, Av \rangle > 0, 
            \end{equation*}
            because $Av \neq O$ since $\langle Av, v \rangle > 0$. 

            Since ${^{t\!}A}^{-1} = A^{-1}$, it follows that $A^{-1}$ is symmetric. 
            Since $A$ is invertible, a given $v$ can be written $v = Aw$ for some $w$ (namely $w = A^{-1}v$). 
            Then 
            \begin{equation*}
                \langle A^{-1}v, v \rangle = \langle A^{-1}Aw, Aw \rangle = \langle w, Aw \rangle = \langle Aw, w \rangle > 0. 
            \end{equation*}
            Hence $A^{-1}$ is positive definite. 
        \end{proof}
    \end{Solution}
    \item 
    Let $U: \mathbf{R}^n \to \mathbf{R}^n$ be a linear map, and let $\langle \ , \ \rangle$ denote the usual scalar (dot) product. 
    Show that the following conditions are equivalent: 
    \er{
        \item 
        $\norm{Uv} = \norm{v}$ for all $v \in \mathbf{R}^n$. 
        \item 
        $\langle Uv, Uw \rangle = \langle v, w \rangle$ for all $v, w \in \mathbf{R}^n$. 
        \item 
        $U$ is invertible, and ${^{t\!}U} = U^{-1}$. 
    }
    [\textit{Hint}: For (ii), use the identity 
    \begin{equation*}
        \langle (v + w), (v + w) \rangle - \langle (v - w), (v - w) \rangle = 4\langle v, w \rangle, 
    \end{equation*}
    and similarly with a $U$ in front of each vector.] 
    When $U$ satisfies any one (and hence all) of these conditions, then $U$ is called \textbf{unitary}. 
    The first condition says that $U$ preserves the norm, and the second says that $U$ preserves the scalar product. 
    \begin{Solution}
        \begin{proof}
            Assume (i).             
            From the identity in the hint, we get 
            \begin{equation*}
                \begin{aligned}
                    4 \langle Uv, Uw \rangle &= \langle U(v + w), U(v + w) \rangle - \langle U(v - w), U(v - w) \rangle \\
                    &= \norm{U(v + w)}^2 - \norm{U(v - w)}^2 \\
                    &= \norm{v + w}^2 - \norm{v - w}^2 \\
                    &= \langle (v + w), (v + w) \rangle - \langle (v - w), (v - w) \rangle \\
                    &= 4\langle v, w \rangle. 
                \end{aligned}
            \end{equation*}
            Hence $\langle Uv, Uw \rangle = \langle v, w \rangle$. 
            The converse is immediate. 

            Assume (i). 
            Then $\Ker U = O$ because if $Uv = O$ then $\norm{v} = 0$ so $v = O$. 
            But a linear map with $O$ kernel from a finite dimensional vector space into itself is invertible, so $U$ is invertible. 
            Also, for all $v, w \in V$, 
            \begin{equation*}
                \langle Uv, Uw \rangle = \langle v, {^{t\!}U}Uw \rangle \quad \text{and also equals} \quad \langle v, w \rangle \quad \text{by hypothesis}. 
            \end{equation*}
            Hence ${^{t\!}U}U = I$ so ${^{t\!}U} = U^{-1}$. Conversely, from ${^{t\!}U} = U^{-1}$ we get 
            \begin{equation*}
                \norm{Uv} = \langle Uv, Uv \rangle = \langle {^{t\!}U}Uv, v \rangle = \langle v, v \rangle = \norm{v}, 
            \end{equation*}
            so $U$ satrisfies (i). 
        \end{proof}
    \end{Solution}
    \item 
    Let $A: \mathbf{R}^n \to \mathbf{R}^n$ be an invertible linear map. 
    \er{
        \item 
        Show that ${^{t\!}A}A$ is symmetric positive definite. 
        \item 
        By Exercise 3b, there is a symmetric positive definite $B$ such that $B^2 = {^{t\!}A}A$. 
        Let $U = AB^{-1}$. 
        Show that $U$ is unitary. 
        \item 
        Show that $A = UB$. 
    }
    \begin{Solution}
        \er{
            \item 
            \begin{proof}
                Since ${^{t\!}({^{t\!}A}A}) = {^{t\!}A}{^{t\!}{^{t\!}A}} = {^{t\!}A}A$, it follows that ${^{t\!}A}A$ is symmetric. 
                Furthermore, for $v \neq O$, we have 
                \begin{equation*}
                    \langle {^{t\!}A}Av, v \rangle = \langle Av, Av \rangle > 0 
                \end{equation*}
                because $A$ is invertible, $Av \neq O$. 
                Hence ${^{t\!}A}A$ is positive definite. 
            \end{proof}
            \item 
            \begin{proof}
                Let $U = AB^{-1}$ where $B^2 = {^{t\!}A}A$. 
                Then 
                \begin{equation*}
                    \begin{aligned}
                        \langle Uv, Uv \rangle &= \langle AB^{-1}v, AB^{-1}v \rangle \\
                        &= \langle B^{-1}v, {^{t\!}A}AB^{-1}v \rangle \\
                        &= \langle B^{-1}v, Bv \rangle \\
                        &= \langle {^{t\!}B}B^{-1}v, v \rangle \\
                        &= \langle v, v \rangle, 
                    \end{aligned}
                \end{equation*}
                because $B$ is symmetric. 
                Hence $U$ is unitary. 
            \end{proof}
            \item 
            \begin{proof}
                Since $U = AB^{-1}$, we have 
                \begin{equation*}
                    UB = AB^{-1}B 
                \end{equation*}
                and hence $A = UB$. 
            \end{proof}
        }
    \end{Solution}
    \item 
    Let $B$ be symmetric positive definite and also unitary. 
    Show that $B = I$. 
    \begin{Solution}
        \begin{proof}
            Let $v$ be a non-zero eigenvector $\lambda > 0$. 
            Then 
            \begin{equation*}
                \langle Bv, Bv \rangle = \lambda^2 \langle v, v \rangle = \langle v, v \rangle 
            \end{equation*}
            because $B$ is unitary. 
            Hence $\lambda^2 = 1$. 
            Hence $\lambda = \pm 1$, and since $\lambda$ is positive, $\lambda = 1$. 
            Since $V$ has a basis consisting of eigenvectors, it follows that $B = I$. 
        \end{proof}
    \end{Solution}
}
}